\documentclass[11pt,a4paper]{article}
\usepackage[margin=25mm]{geometry}
\usepackage[T1]{fontenc}
\usepackage{hyperref}
\begin{document}
\title{A0 knowledge distillation -- Phase 4 CJPE methodology}
\author{CS781 project documentation}
\date{Reconstructed from preserved source and run records}
\maketitle
\section{Problem statement and data}
Phase 4 Court Judgment Prediction predicts binary outcome 0/1 from legal case text. The preserved official inputs are Train/train.json, Train/dev.json and Test/test.json. Test labels are private. Original run records report 994 official dev labels and 1,517 test IDs; macro-F1 is the central metric.
\section{Architecture and mechanics}
The twelve-layer causal A0 student learns CJPE labels and a frozen, larger, own D222 LoRMA teacher's class distribution. Teacher targets are created only for training IDs; temperature-scaled KL is mixed with supervised loss. Data flow: Legal case -> A0 own tokenizer/backbone -> Full A0 student tuning with teacher soft targets -> Token evidence pooling -> Ordered chunk BiGRU -> 2-class logits -> prediction.json.
\section{Training design and checkpoint selection}
Full A0 student tuning with teacher soft targets. Executed v3 runner: `code/v3\_t4x2.ipynb`, `code/runner/phase4\_cjpe/kaggle\_runner.py` with `published\_v3\_schedule.json`; the six jobs share `EvidenceCaseModel` in `code/runner/phase4\_cjpe/models.py`.
\section{Exact run configuration and resource tradeoff}
Experiment family a0; adaptation method full; original parent checkpoint SHA-256 1e8aca9800b87e952ed82e54307dd874220845d6a8545d2568bb743993cfb9dc. Configuration: 256 tokens/chunk, 8 chunks/case, 8192 source characters considered, maximum 12000 training cases, 256 holdout cases per class. Optimizer configuration: backbone learning rate 1e-05, head learning rate 0.0003, gradient accumulation 8, max epochs 4, planned training hours 2.5. Recorded 700 optimizer updates and 5600 seen training examples; selected step 700. Selected checkpoint SHA-256 47df27c673d6ee7619853765d7ad3112ece96bfc9bee9780f30dc606bce3de98; parameters 111113098, trainable 111113098. Original run completed after nominal deadline: True.
\section{Measured results}
Original Kaggle run: complete; 700 optimizer updates; selected step 700. Official 994-case dev macro-F1 **0.563179** at recorded threshold 0.317929; default 0.5 dev macro-F1 0.367628. Verified prediction ZIP SHA-256 `4e26e303cdfc2681221be1e6fbf54be9fa07d114e4386f81cd89c128cf686ec2`, exactly 1,517 IDs. The teacher record and training-only target cache are copied in `evidence/`. User-provided Codabench screenshot: **0.5691** (finished submission).
\section{Evaluation mechanics}
Official labeled dev cases 994; class supports [497, 497]; confusion matrix rows true 0/1, columns predicted 0/1: [[384, 113], [305, 192]]. At selected threshold 0.317929, macro-F1 0.563179; at threshold 0.5, macro-F1 0.367628. The report keeps these thresholds separate because threshold choice changes the binary predictions. The submitted ZIP has 1,517 unique test IDs and binary labels; no test labels were available to this run.
\section{Prediction contract and verification}
The original prediction ZIP contains one root-level prediction.json with 1,517 unique official test IDs and integer labels restricted to 0 and 1. This package checks the ZIP CRC and SHA-256 against the archived run/verification record. Full model checkpoint tensors remain outside GitHub; their downloaded backup receipts and hashes live in the local Phase 4 study archives.
\section{Reproducibility and limits}
The local dev score uses a threshold recorded in the run; inspect its calibration/selection record before treating it as independent model-selection evidence. A prepared model notebook is source scaffolding, while `v3\_t4x2.ipynb` and the Kaggle run evidence identify the executed v3 lineage. Reports and companion notebook were created after training; original scores come from copied artifacts. Model weights and course data are preserved locally, not uploaded to GitHub.
\end{document}
